% !TeX root = ../../tesis.tex
% =============================================================================
%  5.2.1.2 NIVEL DE COBERTURA ESPACIAL
% =============================================================================
\subsection{\bajoRevision{Nivel de Cobertura Espacial ($C$)}}
\label{subsec:cobertura-espacial}
\notaRevision{Los resultados de esta subsección están sujetos a revisión con el
tutor. Los valores presentados corresponden a la corrida del modelo del
18 de abril de 2026 con un isócrono de 800~m y pueden ajustarse tras
validación metodológica.}

El \textbf{Nivel de Cobertura Espacial} ($C$) evalúa la equidad territorial de
la red de transporte masivo al cuantificar qué proporción del área de la
demarcación se encuentra dentro de un radio de caminata aceptable respecto
a alguna estación. Se define mediante la siguiente expresión:

\begin{equation}
    C = \frac{A_{\text{cubierta}}}{A_{\text{total}}} \times 100
    \label{eq:cobertura-res}
\end{equation}

donde $A_{\text{cubierta}}$ es el área resultante de la unión de los
isócronos peatonales (\termino{buffers} de 800~m) de todas las estaciones del sistema,
y $A_{\text{total}}$ es la superficie geodésica total de la demarcación
político-administrativa analizada.

El Cuadro~\ref{tab:peores-cdmx} identifica las cinco alcaldías de la
\textsc{cdmx} con menor cobertura del sistema de transporte masivo,
evidenciando las zonas de mayor exclusión territorial.

\begin{table}[H]
    \centering
    \caption[Alcaldías con menor cobertura de transporte masivo (CDMX)]{
    Cinco alcaldías de la Ciudad de México con menor porcentaje de cobertura
    territorial del sistema de transporte masivo. Isócrono: 800~m.
    Corrida: 2026-04-18.}
    \label{tab:peores-cdmx}
    \begin{tabular}{|l|r|r|r|}
        \hline
        \textbf{Demarcación} & \textbf{Área total (km²)} &
        \textbf{Área cubierta (km²)} & \textbf{Cobertura (\%)} \\
        \hline
        Milpa Alta              & 298.01 &  35.0    &  11.75 \\
        La Magdalena Contreras  &  63.37 &  16.33   &  25.76 \\
        Tlalpan                 & 314.25 &  94.31   &  30.01 \\
        Tláhuac                 &  85.78 &  43.85   &  51.12 \\
        Cuajimalpa de Morelos   &  71.10 &  37.12   &  52.20 \\
        \hline
    \end{tabular}
    \fuentefigura{Elaboración propia con \texttt{VFTModel} y datos de
    \texttt{Apimetro}.}
\end{table}

\begin{table}[H]
    \centering
    \caption[Municipios del EDOMEX con menor cobertura de transporte masivo]{
    Cinco municipios del Estado de México con menor porcentaje de cobertura
    territorial del sistema de transporte masivo. Isócrono: 800~m.
    Corrida: 2026-04-18.}
    \label{tab:peores-edomex}
    \begin{tabular}{|l|r|r|r|}
        \hline
        \textbf{Demarcación} & \textbf{Área total (km²)} &
        \textbf{Área cubierta (km²)} & \textbf{Cobertura (\%)} \\
        \hline
        Tonanitla       &   8.52 & 0.00 & 0.00 \\
        Capulhuac       &  21.50 & 0.00 & 0.00 \\
        Calimaya        & 104.26 & 0.00 & 0.00 \\
        Donato Guerra   & 181.36 & 0.00 & 0.00 \\
        Ayapango        &  50.83 & 0.00 & 0.00 \\
        \hline
    \end{tabular}
    \fuentefigura{Elaboración propia con \texttt{VFTModel} y datos de
    \texttt{Apimetro}.}
\end{table}

% =============================================================================
%  Desglose de Cobertura por Sistema de Transporte
% =============================================================================

\subsubsection{Sistema de Transporte Colectivo Metro (METRO)}
\label{subsubsec:cobertura-metro}

El Sistema de Transporte Colectivo Metro constituye la columna vertebral de la
movilidad masiva en la Ciudad de México. Debido a su naturaleza de red ferroviaria
subterránea y elevada con derecho de vía exclusivo, sus estaciones actúan como
nodos estructurales de alta capacidad que articulan los flujos intermodales de la
metrópoli. El análisis de su cobertura espacial mediante isócronos peatonales de
800~m permite identificar las demarcaciones que se benefician directamente de este
sistema, así como aquellas periferias y zonas centrales que presentan una
desconexión crítica respecto a la red de transporte más masiva de la entidad.

El Cuadro~\ref{tab:cobertura-metro-cdmx} detalla las cinco alcaldías de la Ciudad de
México con los niveles más altos de cobertura territorial respecto a la red del
Metro, evidenciando la concentración de la infraestructura ferroviaria en el núcleo
urbano denso de la cuenca.

\begin{table}[H]
    \centering
    \caption[Alcaldías con mayor cobertura del STC Metro (CDMX)]{
    Cinco alcaldías de la Ciudad de México con mayor porcentaje de cobertura
    territorial para el Sistema de Transporte Colectivo Metro. Isócrono: 800~m.
    Corrida: 2026-04-18.}
    \label{tab:cobertura-metro-cdmx}
    \begin{tabular}{|l|r|r|r|}
        \hline
        \textbf{Demarcación} & \textbf{Área total (km²)} & \textbf{Área cubierta (km²)} & \textbf{Cobertura (\%)} \\
        \hline
        Cuauhtémoc          & 32.50 &  28.35 &  87.23 \\
        Benito Juárez       & 26.68 &  21.34 &  79.99 \\
        Venustiano Carranza & 33.84 &  24.39 &  72.77 \\
        Iztacalco           & 23.08 &  11.66 &  50.52 \\
        Azcapotzalco        & 33.50 &  13.80 &  41.18 \\
        \hline
    \end{tabular}
    \fuentefigura{Elaboración propia con \texttt{VFTModel} y datos de \texttt{Apimetro}.}
\end{table}

Por su parte, la distribución geográfica de esta cobertura se visualiza en la
Figura~\ref{fig:mapa-cobertura-metro}, donde se aprecia la concentración de la
infraestructura en la zona central de la cuenca y el drástico desvanecimiento del
servicio hacia las demarcaciones del sur y del poniente.

\begin{figure}[H]
    \centering
    \includegraphics[width=0.85\textwidth]{Figures/Cap5/Cobertura/Mapas/NB02_fig02_METRO_cobertura.png}
    \caption[Mapa de cobertura territorial del STC Metro en la CDMX]{Mapa del
    nivel de cobertura territorial del Sistema de Transporte Colectivo Metro
    por demarcación en la Ciudad de México, calculado a partir de \termino{buffers}
    de caminata peatonal de 800~m.}
    \label{fig:mapa-cobertura-metro}
    \fuentefigura{Elaboración propia con \texttt{VFTModel} utilizando visualización geoespacial.}
\end{figure}

\subsubsection{Metrobús (MB)}
\label{subsubsec:cobertura-mb}

El sistema Metrobús, que opera bajo el modelo de Autobuses de Tránsito Rápido
(\termino{Bus Rapid Transit}, BRT), se ha consolidado como la segunda red de
transporte masivo más importante de la Ciudad de México. Dada su mayor flexibilidad
constructiva frente a los sistemas ferroviarios pesados, su trazado ha logrado
expandirse por corredores viales estratégicos, complementando y alimentando a la
red del Metro. No obstante, al evaluar su cobertura mediante isócronos peatonales
de 800~m, persisten notables brechas territoriales, dejando a varias demarcaciones
periféricas con acceso limitado o nulo a esta red troncal.

El Cuadro~\ref{tab:cobertura-mb-cdmx} presenta las cinco alcaldías con mayor
proporción de su territorio cubierto por el Metrobús, destacando la concentración
del servicio en las zonas centro, norte y oriente de la ciudad, siguiendo los ejes
viales primarios.

\begin{table}[H]
    \centering
    \caption[Alcaldías con mayor cobertura de Metrobús (CDMX)]{
    Cinco alcaldías de la Ciudad de México con mayor porcentaje de cobertura
    territorial para el sistema Metrobús (MB). Isócrono: 800~m.
    Corrida: 2026-04-18.}
    \label{tab:cobertura-mb-cdmx}
    \begin{tabular}{|l|r|r|r|}
        \hline
        \textbf{Demarcación} & \textbf{Área total (km²)} & \textbf{Área cubierta (km²)} & \textbf{Cobertura (\%)} \\
        \hline
        Cuauhtémoc          & 32.50 &  26.74 &  82.27 \\
        Benito Juárez       & 26.68 &  16.21 &  60.74 \\
        Iztacalco           & 23.08 &  13.95 &  60.44 \\
        Venustiano Carranza & 33.84 &  16.10 &  47.58 \\
        Gustavo A. Madero   & 87.84 &  38.75 &  44.12 \\
        \hline
    \end{tabular}
    \fuentefigura{Elaboración propia con \texttt{VFTModel} y datos de \texttt{Apimetro}.}
\end{table}

Visualmente, el alcance espacial del sistema se ilustra en la
Figura~\ref{fig:mapa-cobertura-mb}. El mapa confirma cómo las líneas de BRT se
concentran fuertemente en las zonas centro, norte y oriente —siguiendo los ejes
viales primarios—, mientras que las barreras topográficas y de traza urbana
mantienen a las demarcaciones del surponiente con una cobertura residual o
inexistente.

\begin{figure}[H]
    \centering
    \includegraphics[width=0.85\textwidth]{Figures/Cap5/Cobertura/Mapas/NB02_fig02_MB_cobertura.png}
    \caption[Mapa de cobertura territorial del Metrobús en la CDMX]{Mapa del
    nivel de cobertura territorial del sistema Metrobús por demarcación en la
    Ciudad de México, calculado a partir de \termino{buffers} de caminata
    peatonal de 800~m.}
    \label{fig:mapa-cobertura-mb}
    \fuentefigura{Elaboración propia con \texttt{VFTModel} utilizando visualización geoespacial.}
\end{figure}

\subsubsection{Cablebús (CBB)}
\label{subsubsec:cobertura-cbb}

El Cablebús representa un paradigma distinto dentro de la red de transporte masivo.
Al ser un sistema de teleférico urbano, su diseño no busca abarcar grandes
extensiones del territorio mediante una red interconectada, sino resolver problemas
de accesibilidad muy puntuales en zonas de alta marginación y topografía compleja
(principalmente en la periferia). Por lo tanto, su nivel de cobertura espacial a
escala metropolitana es inherentemente focalizado, actuando como alimentador de
alta capacidad hacia las líneas del Metro.

El Cuadro~\ref{tab:cobertura-cbb-cdmx} presenta las cinco alcaldías con mayor
cobertura del Cablebús. Como es de esperarse por la naturaleza del sistema, la gran
mayoría de las demarcaciones de la Ciudad de México carecen por completo de esta
infraestructura, por lo que incluso los valores más altos son reducidos en términos
porcentuales absolutos.

\begin{table}[H]
    \centering
    \caption[Alcaldías con mayor cobertura de Cablebús (CDMX)]{
    Cinco alcaldías de la Ciudad de México con mayor porcentaje de cobertura
    territorial para el sistema Cablebús (CBB). Isócrono: 800~m.
    Corrida: 2026-04-18.}
    \label{tab:cobertura-cbb-cdmx}
    \begin{tabular}{|l|r|r|r|}
        \hline
        \textbf{Demarcación} & \textbf{Área total (km²)} & \textbf{Área cubierta (km²)} & \textbf{Cobertura (\%)} \\
        \hline
        Gustavo A. Madero   &  87.84 & 11.31 &  12.88 \\
        Miguel Hidalgo      &  46.39 &  5.50 &  11.85 \\
        Iztapalapa          & 113.07 & 12.74 &  11.27 \\
        Álvaro Obregón      &  95.82 &  4.38 &   4.57 \\
        Venustiano Carranza &  33.84 &  0.00 &   0.00 \\
        \hline
    \end{tabular}
    \fuentefigura{Elaboración propia con \texttt{VFTModel} y datos de \texttt{Apimetro}.}
\end{table}

La Figura~\ref{fig:mapa-cobertura-cbb} ilustra esta focalización espacial. El mapa
revela cómo las áreas de influencia del Cablebús se restringen a corredores muy
específicos (como en Gustavo A. Madero, Iztapalapa y Álvaro Obregón), contrastando
drásticamente con el vacío de infraestructura aérea en el resto de la cuenca.

\begin{figure}[H]
    \centering
    \includegraphics[width=0.85\textwidth]{Figures/Cap5/Cobertura/Mapas/NB02_fig02_CBB_cobertura.png}
    \caption[Mapa de cobertura territorial del Cablebús en la CDMX]{Mapa del
    nivel de cobertura territorial del sistema Cablebús por demarcación en la
    Ciudad de México, calculado a partir de \termino{buffers} de caminata
    peatonal de 800~m.}
    \label{fig:mapa-cobertura-cbb}
    \fuentefigura{Elaboración propia con \texttt{VFTModel} utilizando visualización geoespacial.}
\end{figure}

\subsubsection{Red de Transporte de Pasajeros (RTP)}
\label{subsubsec:cobertura-rtp}

La Red de Transporte de Pasajeros (\textsc{rtp}) desempeña un rol fundamental en
la equidad espacial del sistema de movilidad de la Ciudad de México. A diferencia
de los sistemas confinados, la \textsc{rtp} opera como un servicio de superficie
con vocación social y de alimentación, diseñado específicamente para penetrar en
zonas periféricas, de topografía compleja y de menor densidad, donde la
implementación de infraestructura pesada es inviable. Su extensión territorial es
notablemente mayor que la de cualquier otro sistema de la metrópoli.

El Cuadro~\ref{tab:cobertura-rtp-cdmx} presenta las cinco alcaldías con mayor
proporción de cobertura por parte de la \textsc{rtp}, confirmando la concentración
del servicio en el centro y norte de la ciudad, donde la densidad de rutas y paradas
alcanza sus niveles más altos.

\begin{table}[H]
    \centering
    \caption[Alcaldías con mayor cobertura de RTP (CDMX)]{
    Cinco alcaldías de la Ciudad de México con mayor porcentaje de cobertura
    territorial para la Red de Transporte de Pasajeros (RTP). Isócrono: 800~m.
    Corrida: 2026-04-18.}
    \label{tab:cobertura-rtp-cdmx}
    \begin{tabular}{|l|r|r|r|}
        \hline
        \textbf{Demarcación} & \textbf{Área total (km²)} & \textbf{Área cubierta (km²)} & \textbf{Cobertura (\%)} \\
        \hline
        Miguel Hidalgo      &  46.39 &  41.90 &  90.32 \\
        Cuauhtémoc          &  32.50 &  29.32 &  90.22 \\
        Azcapotzalco        &  33.50 &  30.08 &  89.81 \\
        Gustavo A. Madero   &  87.84 &  75.86 &  86.37 \\
        Iztapalapa          & 113.07 &  93.66 &  82.83 \\
        \hline
    \end{tabular}
    \fuentefigura{Elaboración propia con \texttt{VFTModel} y datos de \texttt{Apimetro}.}
\end{table}

La Figura~\ref{fig:mapa-cobertura-rtp} muestra la dispersión territorial de este
sistema. A diferencia del Metro o del Metrobús, el mapa de la \textsc{rtp} revela
una capilaridad mucho mayor hacia el sur y poniente de la ciudad, confirmando su
rol estratégico como conector de las zonas más alejadas de la cuenca con la red
troncal.

\begin{figure}[H]
    \centering
    \includegraphics[width=0.85\textwidth]{Figures/Cap5/Cobertura/Mapas/NB02_fig02_RTP_cobertura.png}
    \caption[Mapa de cobertura territorial de la RTP en la CDMX]{Mapa del
    nivel de cobertura territorial de la Red de Transporte de Pasajeros
    (\textsc{rtp}) por demarcación en la Ciudad de México, calculado a partir
    de \termino{buffers} de caminata peatonal de 800~m.}
    \label{fig:mapa-cobertura-rtp}
    \fuentefigura{Elaboración propia con \texttt{VFTModel} utilizando visualización geoespacial.}
\end{figure}

\subsubsection{Trolebús (TROLE)}
\label{subsubsec:cobertura-trole}

El Servicio de Transportes Eléctricos (\textsc{ste}), a través de su red de
Trolebús, constituye uno de los ejes históricos y fundamentales de la
electromovilidad de cero emisiones en la capital. Su red se distribuye
estratégicamente sobre los principales ejes viales, concentrando su mayor área
de influencia en las zonas centro, norte y oriente de la ciudad.

\textbf{Nota metodológica:} Es importante precisar que, al momento de la
ejecución de este análisis espacial (18 de abril de 2026), el Portal de Datos
Abiertos de la Ciudad de México aún no había liberado los trazados geoespaciales
ni el padrón de paradas oficiales correspondientes a la \textbf{Línea 13} del
Trolebús. Por consiguiente, dicha línea no se encuentra contabilizada en los
isócronos de cobertura presentados en esta sección.

El Cuadro~\ref{tab:cobertura-trole-cdmx} enumera las cinco alcaldías con mayor
cobertura territorial de este sistema, todas ellas ubicadas en la zona central y
oriente de la ciudad, donde se concentra la red eléctrica de superficie.

\begin{table}[H]
    \centering
    \caption[Alcaldías con mayor cobertura de Trolebús (CDMX)]{
    Cinco alcaldías de la Ciudad de México con mayor porcentaje de cobertura
    territorial para el sistema Trolebús (TROLE). Isócrono: 800~m.
    Corrida: 2026-04-18.}
    \label{tab:cobertura-trole-cdmx}
    \begin{tabular}{|l|r|r|r|}
        \hline
        \textbf{Demarcación} & \textbf{Área total (km²)} & \textbf{Área cubierta (km²)} & \textbf{Cobertura (\%)} \\
        \hline
        Cuauhtémoc    &  32.50 & 22.72 &  69.90 \\
        Benito Juárez &  26.68 & 16.50 &  61.85 \\
        Azcapotzalco  &  33.50 & 17.48 &  52.18 \\
        Iztacalco     &  23.08 & 11.93 &  51.69 \\
        Coyoacán      & 125.17 & 25.17 &  46.71 \\
        \hline
    \end{tabular}
    \fuentefigura{Elaboración propia con \texttt{VFTModel} y datos de \texttt{Apimetro}.}
\end{table}

Por su parte, la Figura~\ref{fig:mapa-cobertura-trole} ilustra la huella espacial
del Trolebús. El mapa visibiliza la fuerte concentración del servicio en alcaldías
como Benito Juárez, Cuauhtémoc, Iztacalco e Iztapalapa (impulsada por el Trolebús
Elevado), contrastando drásticamente con la carencia de infraestructura eléctrica
hacia las periferias montañosas.

\begin{figure}[H]
    \centering
    \includegraphics[width=0.85\textwidth]{Figures/Cap5/Cobertura/Mapas/NB02_fig02_TROLE_cobertura.png}
    \caption[Mapa de cobertura territorial del Trolebús en la CDMX]{Mapa del
    nivel de cobertura territorial del sistema Trolebús por demarcación en la
    Ciudad de México, calculado a partir de \termino{buffers} de caminata
    peatonal de 800~m.}
    \label{fig:mapa-cobertura-trole}
    \fuentefigura{Elaboración propia con \texttt{VFTModel} utilizando visualización geoespacial.}
\end{figure}

\subsubsection{Corredores Concesionados (CC)}
\label{subsubsec:cobertura-cc}

Los Corredores Concesionados constituyen la red de rutas de transporte de
superficie operadas por empresas privadas con concesión del Gobierno de la Ciudad
de México. Con base en su lógica de trazado sobre los ejes viales primarios y
secundarios de la ciudad consolidada, este sistema presenta uno de los niveles de
cobertura territorial más altos de la metrópoli, especialmente en las alcaldías de
mayor densidad urbana del centro y el oriente. Más que un sistema residual, los
Corredores Concesionados actúan como la red capilar que penetra en los intersticios
que los sistemas troncales no pueden servir directamente.

El Cuadro~\ref{tab:cobertura-cc-cdmx} presenta las cinco alcaldías con mayor
cobertura territorial de este sistema, donde se aprecia la saturación práctica del
servicio en el núcleo urbano denso de la capital.

\begin{table}[H]
    \centering
    \caption[Alcaldías con mayor cobertura de Corredores Concesionados (CDMX)]{
    Cinco alcaldías de la Ciudad de México con mayor porcentaje de cobertura
    territorial para el sistema de Corredores Concesionados (CC). Isócrono: 800~m.
    Corrida: 2026-04-18.}
    \label{tab:cobertura-cc-cdmx}
    \begin{tabular}{|l|r|r|r|}
        \hline
        \textbf{Demarcación} & \textbf{Área total (km²)} & \textbf{Área cubierta (km²)} & \textbf{Cobertura (\%)} \\
        \hline
        Cuauhtémoc          & 32.50 & 32.47 &  99.90 \\
        Benito Juárez       & 26.68 & 26.32 &  98.65 \\
        Iztacalco           & 23.08 & 22.25 &  96.42 \\
        Miguel Hidalgo      & 46.39 & 44.33 &  95.57 \\
        Venustiano Carranza & 33.84 & 27.77 &  82.06 \\
        \hline
    \end{tabular}
    \fuentefigura{Elaboración propia con \texttt{VFTModel} y datos de \texttt{Apimetro}.}
\end{table}

La Figura~\ref{fig:mapa-cobertura-cc} confirma la distribución omnipresente de
este sistema en la ciudad consolidada. El mapa ilustra cómo los Corredores
Concesionados recubren de manera prácticamente total las alcaldías centrales,
mientras que el servicio se va diluyendo hacia las periferias con suelo de
conservación o menor densidad.

\begin{figure}[H]
    \centering
    \includegraphics[width=0.85\textwidth]{Figures/Cap5/Cobertura/Mapas/NB02_fig02_CC_cobertura.png}
    \caption[Mapa de cobertura territorial de Corredores Concesionados en la CDMX]{
    Mapa del nivel de cobertura territorial del sistema de Corredores
    Concesionados por demarcación en la Ciudad de México, calculado a partir de
    \termino{buffers} de caminata peatonal de 800~m.}
    \label{fig:mapa-cobertura-cc}
    \fuentefigura{Elaboración propia con \texttt{VFTModel} utilizando visualización geoespacial.}
\end{figure}

\subsubsection{Tren Ligero (TL)}
\label{subsubsec:cobertura-tl}

El Tren Ligero es un sistema de transporte férreo de capacidad media operado por
el Servicio de Transportes Eléctricos (\textsc{ste}). Su trazado consta de una
única línea que conecta la estación terminal Tasqueña de la Línea 2 del Metro con
el centro de la alcaldía Xochimilco. Debido a esta configuración lineal y su
confinamiento exclusivo en el sur de la cuenca, su área de influencia geográfica
es altamente focalizada.

El Cuadro~\ref{tab:cobertura-tl-cdmx} presenta las alcaldías con mayor nivel de
cobertura por parte del Tren Ligero. Únicamente Coyoacán, Xochimilco y Tlalpan
registran acceso al sistema dentro del isócrono peatonal de 800~m, mientras que
el resto de las demarcaciones carece por completo de este servicio.

\begin{table}[H]
    \centering
    \caption[Alcaldías con mayor cobertura de Tren Ligero (CDMX)]{
    Cinco alcaldías de la Ciudad de México con mayor porcentaje de cobertura
    territorial para el sistema Tren Ligero (TL). Isócrono: 800~m.
    Corrida: 2026-04-18.}
    \label{tab:cobertura-tl-cdmx}
    \begin{tabular}{|l|r|r|r|}
        \hline
        \textbf{Demarcación} & \textbf{Área total (km²)} & \textbf{Área cubierta (km²)} & \textbf{Cobertura (\%)} \\
        \hline
        Coyoacán            &  53.88 & 9.31 &  17.28 \\
        Xochimilco          & 114.03 & 6.82 &   5.98 \\
        Tlalpan             & 314.25 & 4.07 &   1.30 \\
        Venustiano Carranza &  33.84 & 0.00 &   0.00 \\
        Tláhuac             &  85.78 & 0.00 &   0.00 \\
        \hline
    \end{tabular}
    \fuentefigura{Elaboración propia con \texttt{VFTModel} y datos de \texttt{Apimetro}.}
\end{table}

La Figura~\ref{fig:mapa-cobertura-tl} ilustra espacialmente esta restricción
territorial. El mapa confirma cómo el servicio es exclusivo del corredor sur,
actuando principalmente como alimentador de la red masiva central para los
habitantes de las alcaldías sureñas.

\begin{figure}[H]
    \centering
    \includegraphics[width=0.85\textwidth]{Figures/Cap5/Cobertura/Mapas/NB02_fig02_TL_cobertura.png}
    \caption[Mapa de cobertura territorial del Tren Ligero en la CDMX]{Mapa del
    nivel de cobertura territorial del sistema Tren Ligero por demarcación en
    la Ciudad de México, calculado a partir de \termino{buffers} de caminata
    peatonal de 800~m.}
    \label{fig:mapa-cobertura-tl}
    \fuentefigura{Elaboración propia con \texttt{VFTModel} utilizando visualización geoespacial.}
\end{figure}

\subsubsection{Mexibús (MXB)}
\label{subsubsec:cobertura-mxb}

El sistema Mexibús es el equivalente al Metrobús en el Estado de México: una red
de autobuses de tránsito rápido (\termino{Bus Rapid Transit}, BRT) que opera sobre
corredores viales del ámbito metropolitano conurbado. Su presencia dentro de los
límites político-administrativos de la Ciudad de México es marginal, circunscrita
únicamente a los puntos de cruce o de conexión interestatal, por lo que su
cobertura territorial en las alcaldías de la capital es forzosamente residual.

El Cuadro~\ref{tab:cobertura-mxb-cdmx} identifica las alcaldías de la Ciudad de
México con mayor cobertura territorial del sistema Mexibús. Los valores reducidos
—ninguno superior al 7\%— confirman que este sistema opera fundamentalmente fuera
de los límites de la entidad y que su incidencia en la cobertura intraurbana
capitalina es mínima.

\begin{table}[H]
    \centering
    \caption[Alcaldías con mayor cobertura de Mexibús (CDMX)]{
    Cinco alcaldías de la Ciudad de México con mayor porcentaje de cobertura
    territorial para el sistema Mexibús (MXB). Isócrono: 800~m.
    Corrida: 2026-04-18.}
    \label{tab:cobertura-mxb-cdmx}
    \begin{tabular}{|l|r|r|r|}
        \hline
        \textbf{Demarcación} & \textbf{Área total (km²)} & \textbf{Área cubierta (km²)} & \textbf{Cobertura (\%)} \\
        \hline
        Iztacalco           &  23.08 & 1.52 & 6.58 \\
        Venustiano Carranza &  33.84 & 1.90 & 5.62 \\
        Gustavo A. Madero   &  87.84 & 3.23 & 3.67 \\
        Xochimilco          & 114.03 & 0.00 & 0.00 \\
        Tláhuac             &  85.78 & 0.00 & 0.00 \\
        \hline
    \end{tabular}
    \fuentefigura{Elaboración propia con \texttt{VFTModel} y datos de \texttt{Apimetro}.}
\end{table}

La Figura~\ref{fig:mapa-cobertura-mxb} confirma esta presencia puntual del
sistema dentro de la capital. El mapa ilustra cómo la única área de influencia
del Mexibús en la Ciudad de México corresponde a los cruces fronterizos del
nororiente, sin penetrar en el tejido urbano interior.

\begin{figure}[H]
    \centering
    \includegraphics[width=0.85\textwidth]{Figures/Cap5/Cobertura/Mapas/NB02_fig02_MEXIBÚS_cobertura.png}
    \caption[Mapa de cobertura territorial de Mexibús en la CDMX]{Mapa del
    nivel de cobertura territorial del sistema Mexibús por demarcación en la
    Ciudad de México, calculado a partir de \termino{buffers} de caminata
    peatonal de 800~m.}
    \label{fig:mapa-cobertura-mxb}
    \fuentefigura{Elaboración propia con \texttt{VFTModel} utilizando visualización geoespacial.}
\end{figure}

\subsubsection{Mexicable (MEXICABLE)}
\label{subsubsec:cobertura-mexicable}

El Mexicable es un sistema de teleférico urbano que opera principalmente en el
Estado de México, diseñado para sortear las barreras topográficas de las zonas
altas y densamente pobladas de la periferia metropolitana. Su incursión en la
Ciudad de México es mínima y de carácter estrictamente interconectivo, fungiendo
como enlace alimentador hacia las redes troncales en la zona norte de la capital.

El Cuadro~\ref{tab:cobertura-mexicable-cdmx} presenta las alcaldías de la capital
con mayor nivel de cobertura para el Mexicable. Como es natural dada su jurisdicción
y trazado, únicamente la alcaldía Gustavo A. Madero registra un área de influencia
efectiva dentro del isócrono peatonal de 800~m, mientras que el resto de las
demarcaciones presentan cobertura nula.

\begin{table}[H]
    \centering
    \caption[Alcaldías con mayor cobertura de Mexicable (CDMX)]{
    Cinco alcaldías de la Ciudad de México con mayor porcentaje de cobertura
    territorial para el sistema Mexicable. Isócrono: 800~m.
    Corrida: 2026-04-18.}
    \label{tab:cobertura-mexicable-cdmx}
    \begin{tabular}{|l|r|r|r|}
        \hline
        \textbf{Demarcación} & \textbf{Área total (km²)} & \textbf{Área cubierta (km²)} & \textbf{Cobertura (\%)} \\
        \hline
        Gustavo A. Madero   &  87.84 & 2.80 & 3.19 \\
        Venustiano Carranza &  33.84 & 0.00 & 0.00 \\
        Xochimilco          & 114.03 & 0.00 & 0.00 \\
        Tláhuac             &  85.78 & 0.00 & 0.00 \\
        Tlalpan             & 314.25 & 0.00 & 0.00 \\
        \hline
    \end{tabular}
    \fuentefigura{Elaboración propia con \texttt{VFTModel} y datos de \texttt{Apimetro}.}
\end{table}

La Figura~\ref{fig:mapa-cobertura-mexicable} visualiza la focalización extrema de
este servicio. El mapa ilustra claramente cómo la infraestructura aérea del
Mexicable se restringe a un pequeño polígono en el extremo norte de la ciudad,
reiterando su naturaleza como transporte de penetración desde la zona conurbada.

\begin{figure}[H]
    \centering
    \includegraphics[width=0.85\textwidth]{Figures/Cap5/Cobertura/Mapas/NB02_fig02_MEXICABLE_cobertura.png}
    \caption[Mapa de cobertura territorial del Mexicable en la CDMX]{Mapa del
    nivel de cobertura territorial del sistema Mexicable por demarcación en la
    Ciudad de México, calculado a partir de \termino{buffers} de caminata
    peatonal de 800~m.}
    \label{fig:mapa-cobertura-mexicable}
    \fuentefigura{Elaboración propia con \texttt{VFTModel} utilizando visualización geoespacial.}
\end{figure}

\subsubsection{Tren Suburbano (SUB)}
\label{subsubsec:cobertura-sub}

El Tren Suburbano de la \zmvm{} es un sistema ferroviario regional de alta
capacidad que conecta la terminal Buenavista, en la alcaldía Cuauhtémoc, con el
municipio de Cuautitlán, en el Estado de México. Su trazado, mayoritariamente
confinado al corredor norte de la cuenca metropolitana, limita su área de
influencia directa dentro de la Ciudad de México a un corredor muy acotado. Para
las alcaldías del sur y oriente, este sistema resulta prácticamente inaccesible
sin una transferencia modal significativa.

El Cuadro~\ref{tab:cobertura-sub-cdmx} detalla las cinco alcaldías de la Ciudad de
México con mayor nivel de cobertura por parte del Tren Suburbano, evidenciando la
concentración casi exclusiva del servicio en las demarcaciones del norte y centro
de la capital.

\begin{table}[H]
    \centering
    \caption[Alcaldías con mayor cobertura de Tren Suburbano (CDMX)]{
    Cinco alcaldías de la Ciudad de México con mayor porcentaje de cobertura
    territorial para el sistema Tren Suburbano (SUB). Isócrono: 800~m.
    Corrida: 2026-04-18.}
    \label{tab:cobertura-sub-cdmx}
    \begin{tabular}{|l|r|r|r|}
        \hline
        \textbf{Demarcación} & \textbf{Área total (km²)} & \textbf{Área cubierta (km²)} & \textbf{Cobertura (\%)} \\
        \hline
        Cuauhtémoc          &  32.50 & 2.01 & 6.18 \\
        Azcapotzalco        &  33.50 & 2.01 & 5.99 \\
        Xochimilco          & 114.03 & 0.00 & 0.00 \\
        Venustiano Carranza &  33.84 & 0.00 & 0.00 \\
        Tláhuac             &  85.78 & 0.00 & 0.00 \\
        \hline
    \end{tabular}
    \fuentefigura{Elaboración propia con \texttt{VFTModel} y datos de \texttt{Apimetro}.}
\end{table}

La Figura~\ref{fig:mapa-cobertura-sub} confirma la huella territorial de este
sistema. El mapa evidencia cómo la cobertura del Tren Suburbano apenas alcanza la
alcaldía Cuauhtémoc y el extremo norte de Azcapotzalco, mientras que el resto del
territorio capitalino permanece fuera de su radio de influencia peatonal.

\begin{figure}[H]
    \centering
    \includegraphics[width=0.85\textwidth]{Figures/Cap5/Cobertura/Mapas/NB02_fig02_SUB_cobertura.png}
    \caption[Mapa de cobertura territorial del Tren Suburbano en la CDMX]{Mapa
    del nivel de cobertura territorial del sistema Tren Suburbano por demarcación
    en la Ciudad de México, calculado a partir de \termino{buffers} de caminata
    peatonal de 800~m.}
    \label{fig:mapa-cobertura-sub}
    \fuentefigura{Elaboración propia con \texttt{VFTModel} utilizando visualización geoespacial.}
\end{figure}

\subsubsection{Tren Interurbano (INTERURBANO)}
\label{subsubsec:cobertura-interurbano}

El Tren Interurbano ``El Insurgente'' es un sistema ferroviario de carácter
regional diseñado para conectar el Valle de Toluca con la Ciudad de México. Su
trazo dentro de la capital es incipiente y se concentra exclusivamente en el sector
poniente, sirviendo como puerta de entrada a la metrópoli. Debido a su vocación
interurbana, su impacto en la cobertura espacial a nivel de alcaldías es altamente
localizado.

El Cuadro~\ref{tab:cobertura-interurbano-cdmx} presenta las alcaldías con mayor
nivel de cobertura para este sistema. Como reflejo de su ruta de ingreso, únicamente
Cuajimalpa de Morelos y Álvaro Obregón presentan áreas de influencia dentro del
isócrono peatonal de 800~m, dejando al resto de la ciudad sin cobertura directa
por este medio.

\begin{table}[H]
    \centering
    \caption[Alcaldías con mayor cobertura de Tren Interurbano (CDMX)]{
    Cinco alcaldías de la Ciudad de México con mayor porcentaje de cobertura
    territorial para el sistema Tren Interurbano (INTERURBANO). Isócrono: 800~m.
    Corrida: 2026-04-18.}
    \label{tab:cobertura-interurbano-cdmx}
    \begin{tabular}{|l|r|r|r|}
        \hline
        \textbf{Demarcación} & \textbf{Área total (km²)} & \textbf{Área cubierta (km²)} & \textbf{Cobertura (\%)} \\
        \hline
        Cuajimalpa de Morelos &  71.10 & 1.59 & 2.23 \\
        Álvaro Obregón        &  95.82 & 0.42 & 0.44 \\
        Venustiano Carranza   &  33.84 & 0.00 & 0.00 \\
        Xochimilco            & 114.03 & 0.00 & 0.00 \\
        Tláhuac               &  85.78 & 0.00 & 0.00 \\
        \hline
    \end{tabular}
    \fuentefigura{Elaboración propia con \texttt{VFTModel} y datos de \texttt{Apimetro}.}
\end{table}

La Figura~\ref{fig:mapa-cobertura-interurbano} visualiza este alcance espacial
limitado. El mapa hace evidente la focalización de la infraestructura en el límite
occidental de la cuenca, evidenciando su papel como colector foráneo más que como
distribuidor urbano capilar.

\begin{figure}[H]
    \centering
    \includegraphics[width=0.85\textwidth]{Figures/Cap5/Cobertura/Mapas/NB02_fig02_INTERURBANO_cobertura.png}
    \caption[Mapa de cobertura territorial del Tren Interurbano en la CDMX]{Mapa
    del nivel de cobertura territorial del sistema Tren Interurbano por
    demarcación en la Ciudad de México, calculado a partir de \termino{buffers}
    de caminata peatonal de 800~m.}
    \label{fig:mapa-cobertura-interurbano}
    \fuentefigura{Elaboración propia con \texttt{VFTModel} utilizando visualización geoespacial.}
\end{figure}

\subsubsection{Pumabús (PB)}
\label{subsubsec:cobertura-pb}

El sistema Pumabús es un servicio de transporte interno operado por la Universidad
Nacional Autónoma de México (UNAM), de carácter público y gratuito, cuya operación
es exclusiva del campus de Ciudad Universitaria. No obstante, debido a su relevancia
modal y estructural en el sur de la ciudad, se encuentra registrado formalmente
dentro del Portal de Datos Abiertos de la Ciudad de México, lo que permite integrar
su infraestructura de paradas y circuitos en el análisis macro de la conectividad
metropolitana.

El Cuadro~\ref{tab:cobertura-pb-cdmx} identifica las cinco alcaldías de la Ciudad
de México con mayor cobertura territorial respecto a la red del Pumabús,
evidenciando su completa focalización en el núcleo universitario y su vecindad
inmediata.

\begin{table}[H]
    \centering
    \caption[Alcaldías con mayor cobertura de Pumabús (CDMX)]{
    Cinco alcaldías de la Ciudad de México con mayor porcentaje de cobertura
    territorial para el sistema Pumabús (PB). Isócrono: 800~m.
    Corrida: 2026-04-18.}
    \label{tab:cobertura-pb-cdmx}
    \begin{tabular}{|l|r|r|r|}
        \hline
        \textbf{Demarcación} & \textbf{Área total (km²)} & \textbf{Área cubierta (km²)} & \textbf{Cobertura (\%)} \\
        \hline
        Coyoacán            &  53.88 & 10.85 &  20.13 \\
        Álvaro Obregón      &  95.82 &  2.25 &   2.34 \\
        Tlalpan             & 314.25 &  0.25 &   0.08 \\
        Venustiano Carranza &  33.84 &  0.00 &   0.00 \\
        Xochimilco          & 114.03 &  0.00 &   0.00 \\
        \hline
    \end{tabular}
    \fuentefigura{Elaboración propia con \texttt{VFTModel} y datos de \texttt{Apimetro}.}
\end{table}

La configuración geoespacial de esta influencia se presenta en la
Figura~\ref{fig:mapa-cobertura-pb}. El mapa hace patente la delimitación del
servicio al polígono escolar de Ciudad Universitaria, extendiendo su área de
influencia peatonal de 800~m únicamente hacia las franjas fronterizas de las tres
alcaldías circundantes.

\begin{figure}[H]
    \centering
    \includegraphics[width=0.85\textwidth]{Figures/Cap5/Cobertura/Mapas/NB02_fig02_PUMABUS_cobertura.png}
    \caption[Mapa de cobertura territorial del Pumabús en la CDMX]{Mapa del
    nivel de cobertura territorial del sistema Pumabús por demarcación en la
    Ciudad de México, calculado a partir de \termino{buffers} de caminata
    peatonal de 800~m.}
    \label{fig:mapa-cobertura-pb}
    \fuentefigura{Elaboración propia con \texttt{VFTModel} utilizando visualización geoespacial.}
\end{figure}
